% Pursuing Stacks (cotes 134-2 to 134-8): archive cote 134-6, site page 37 (PDF page 38), typescript page 355.
% typescript. Draft compared: ps5.tex:3153-3221.
% AI-transcribed by Claude Opus 5.5 (claude-opus-5-5, adaptive thinking, effort high), Message Batches API, batch ps3 (msgbatch_01XdFmuVVBAcgc7EFxdtCqsj), 2026-11-12; not yet checked by eye.
\scanpage{134-6/37}
\mspage{355}

The property \xout{that} for a topos X, with category of sheaves \hand{$\mathcal{A}$}, that
any object F in \hand{$\mathcal{A}$} \xout{such that} covering the final`object e\hand{$_{\mathcal{A}}$} \inserted{should} ha\over{s}{ve} a section,
can be expressed by saying that the latter is a \uline{projective} object in
the category \hand{$\mathcal{A}$}. Following the principle to use the same names for
properties of a topos, and corresponding properties of the final
sheaf on it, we may call a topos with the above property a \uline{projective}
\uline{topos}. Thus the``non-empty'' topo\"{\i} such that every ``non-empty'' sheaf \xout{xx}
has a section, are exactly the topo\"{\i} which are both vertical and projec-%
-tive.

* \quad *

*

Here is the promised ``exact relationship'' between the three
standard homotopy structures h,h',h'' on A\^{}, where A is any small cate-%
gory (cf. end of section 96, p. 354) :

\margin{\'ecarter}
\diagram{\begin{tikzcd}[column sep=large]
(\hand{*}) & h \arrow[d, equal] \arrow[r, phantom, "\xout{x}(A\ \text{tot.}\ \underline{W}\text{-asph.})"{above}, "\hand{$\leq$}"{description}] & h' \arrow[d, equal] \arrow[r, phantom, "(\underline{W}\subset\underline{W}_o)"{above}, "\hand{$\leq$}"{description}] & h'' \arrow[d, equal] \arrow[r, phantom, "\inserted[typed]{(A\ \text{vertical})}\ \xout{(A\ \text{vertical})}"{above}, "\hand{$\leq$}"{description}] & h \\
 & h_{A\hat{}\hand{,}A} & h_{\underline{W}_A} & h_{\underline{W}_{oA}}=h_{A\hat{}} & ,
\end{tikzcd}}

where \uline{W} is a given basic localizer. Above each one of the three \inserted{conditional} ine-%
qualities between homotopy structures h,h',h'' I wrote the natural \xout{assu}
assumption \inserted{on A or \uline{W}} validating it , and \xout{underneath of} in the diagram I have
recalled the definition of the three homotopy structures. A propos
the description $h = h_{A\hat{},A}$ , the notation used here is $h_{M,A}$ when
M is a category stable under finite products and A a full subcategory,
for designating the homotopy structure on M generated by intervals in
M coming from A. A propos $h'' = h_{A\hat{}}$ , I recall the notation $h_M$ for
designating the canonical homotopy structure on a category M satisfying
suitable conditions (section 57). Also I recall
\[ \underline{W}_o = \hand{$\{$} f\hand{$\in$}\mathrm{Fl}(\mathrm{Cat}) \hand{$|$} \hand{$\pi$}_o(f) \text{ bijective}\hand{$\}$} \ . \]
The diagram implies that if $\underline{W}\hand{$\subset$}\underline{W}_o$ i.e. \uline{W} satisfies Loc 4), and if
moreover A is vertical and \xout{totally Wx non empty (hence} tota\unsure{\over{l}{ll}}y $\underline{W}$\xout{$_{oo}$}-as-%
pheric, \inserted{then all three homotopy structures coincide. Also, taking} \deleted{\xout{\unsure{by what we just saw, and a fortiori $\underline{W}$-aspheric. If we admit here}}}
\deleted{\xout{that necessarily \unsure{$\underline{W}$ $\underline{W}_{oo}$}}} $\underline{W} = \underline{W}_o$ , we see that h=h'' if A is verti-%
cal and totally 0-connected, w\unsure{\hand{h}}ich is \unsure{\over{L}{l}}emma 2 of p.268 for A\^{}, A .

From (1) it follows of course that if A is a contractor, then
(for any \uline{W} satisfying Loc 4)) the three homotopy structures h,h',h''
on A\^{} coincide. In case the contractor A is not trivial, hence A is

\margin{37.}
